\section{Introduction}

Whether out of scientific interest, curiosity, or to raise awareness of space-related risks, the observation of near-Earth objects is a topic of growing interest. Near-Earth Objects (NEOs) are asteroids and comets whose orbits can bring them close to Earth. Because of the potential consequences of an impact, their detection and tracking are the subject of constant monitoring, in particular by NASA, which makes a large amount of data on NEOs available to the public through its NeoWs API.

However, raw access to this data is not enough to meet a user's needs. Quickly identifying the next close approaches, searching for a specific object, following its distance to Earth over time or being notified of an event of interest all require a tool that processes the data and presents it in an accessible way. NEO-Watch was designed to fill this gap: it is a web application that imports NASA's data into its own database and lets each user explore, organize and personalize it.

In September, our analysis report defined the scope of the application, the choices imposed by the NASA API and the overall design. This final report presents the application as it was actually built. NEO-Watch covers the six mandatory features (F1 to F6), from account management to personalized alerts, as well as the three optional features we chose: statistics on the NEOs stored in the database (FO3), automatic daily updates of the data (FO7) and the NASA Astronomy Picture of the Day (FO8). It relies on the layered architecture imposed by the project template: a Streamlit frontend, a FastAPI backend organized into controllers, services and DAO, and a PostgreSQL database.

The report is organized as follows. Section 2 recalls the requirements and the user roles, and walks through a complete usage scenario. Section 3 presents the final design of the application: diagrams, architecture and data model. Section 4 details the technical implementation, in particular the import of NASA data and security, and Section 5 describes our testing strategy and code quality tools. Section 6 is devoted to the team's organization and compares the planned schedule with the actual one. Finally, the conclusion reviews the differences between the analysis and the final application, and opens up perspectives for future improvements.
